
\documentclass[12pt,a4paper]{article}

\usepackage[a4paper,margin=1in]{geometry}
\usepackage{graphicx}
\usepackage{booktabs}
\usepackage{amsmath,amssymb}
\usepackage{float}
\usepackage{hyperref}
\usepackage{xcolor}
\usepackage{listings}
\usepackage{enumitem}
\usepackage{fancyhdr}
\usepackage{longtable}

\hypersetup{colorlinks=true,linkcolor=blue,urlcolor=blue}

\pagestyle{fancy}
\fancyhf{}
\lhead{ICS1512}
\rhead{Machine Learning Laboratory}
\cfoot{\thepage}

\lstset{
language=Python,
basicstyle=\ttfamily\small,
numbers=left,
frame=single,
breaklines=true,
showstringspaces=false
}

\begin{document}

\begin{tabular}{|p{4cm}|p{11cm}|}
\hline
Experiment No. & 7 \\ \hline
Experiment Title & Dimensionality Reduction and Model Evaluation (With and Without PCA)
\footnote{\textbf{GitHub Repository: }\url{https://github.com/Barath-j593/ML-assignments}}\\ \hline
Student Name & NavinKumar S \\ \hline
Register Number & 3122247001039 \\ \hline
Faculty & Dr. Poreddy Ajay Kumar Reddy \\ \hline
Submission Date & 19/09/26 \\ \hline
\end{tabular}

\section{Objective}
To study the effect of dimensionality reduction using Principal Component Analysis (PCA) on
the performance of ten machine learning classifiers, by training and validating each model
both without PCA (original feature space) and with PCA (reduced feature space), performing
hyperparameter tuning and 5-fold cross-validation in both settings.

\section{Problem Statement}
Given the Wisconsin Diagnostic Breast Cancer dataset of 30 numerical features, many of which
are highly correlated (mean/worst/standard-error versions of the same 10 underlying cell
measurements), the problem is to determine whether compressing this feature space with PCA
improves, degrades, or leaves unchanged the accuracy, F1-score, and cross-validation
stability of ten different classifiers -- spanning linear, distance-based, probabilistic,
tree-based, and ensemble model families.

\section{Dataset Description}
\begin{longtable}{|p{6cm}|p{8cm}|}
\hline
Dataset Name & Wisconsin Diagnostic Breast Cancer (WDBC) \\ \hline
Dataset Source & UCI Machine Learning Repository -- \url{https://archive.ics.uci.edu} (loaded via \texttt{sklearn.datasets.load\_breast\_cancer}) \\ \hline
Number of Samples & 569 \\ \hline
Number of Features & 30 numerical attributes (mean, standard error, and ``worst'' value of 10 cell-nucleus measurements) \\ \hline
Number of Classes & 2 (357 benign / 212 malignant) \\ \hline
Missing Values & 0 \\ \hline
Train-Test Split & 80:20 stratified (455 train / 114 test), \texttt{random\_state=42} \\ \hline
\end{longtable}

\textbf{Preprocessing steps:} Label encoding was not required (target is already numeric,
0=malignant, 1=benign). No missing values were present. All 30 features were standardized
(\texttt{StandardScaler}, zero mean, unit variance, fit on the training split only) before
both the No-PCA and With-PCA branches, since several of the evaluated models (SVM, KNN,
Logistic Regression) and PCA itself are scale-sensitive.

\textbf{Note on XGBoost:} The training environment for this experiment does not have network
access to install the third-party \texttt{xgboost} package. As a substitute, scikit-learn's
built-in \texttt{HistGradientBoostingClassifier} was used -- a histogram-based gradient
boosting implementation from the same algorithmic family as XGBoost/LightGBM, offering a
comparable accuracy/speed profile. This substitution is noted throughout as ``XGBoost
(HistGB)''.

\section{PCA Design Choice}

\begin{figure}[H]
\centering
\includegraphics[width=0.7\linewidth]{pca_scree_plot.png}
\caption{PCA scree plot -- cumulative explained variance vs. number of components}
\end{figure}

\begin{longtable}{|p{4cm}|p{4cm}|p{3cm}|p{4cm}|}
\hline
\textbf{Setting} & \textbf{Chosen Components / Variance Target} & \textbf{Explained Variance (\%)} & \textbf{Justification} \\ \hline
\endhead
With-PCA & 95\% cumulative variance target $\rightarrow$ 10 components & 95.27\% & A 95\% variance retention threshold is a standard default that keeps almost all discriminative information while still compressing 30 raw features into just 10 principal components (a 3$\times$ reduction) -- justified further by the scree plot's visible ``elbow'' flattening after roughly 6--10 components, beyond which each additional component contributes very little further variance. \\ \hline
\end{longtable}
\textbf{Inference:} 10 of 30 components already capture over 95\% of total variance,
confirming the strong inter-feature correlation observed in the EDA correlation heatmap of
earlier experiments on this dataset (mean/worst/error versions of the same measurement are,
unsurprisingly, highly redundant) -- exactly the situation where PCA is expected to be most
effective at compressing without much information loss.

\section{Exploratory Data Analysis}
\begin{figure}[H]
\centering
\includegraphics[width=0.5\linewidth]{class_distribution.png}
\caption{Class distribution -- Breast Cancer Wisconsin (Diagnostic)}
\end{figure}
\textbf{Inference:} The dataset is mildly imbalanced (357 benign vs. 212 malignant, about
63:37), consistent across all experiments conducted on this dataset. No missing-value
visualization is required since there are zero missing values.

\section{Hyperparameter Grids}
\begin{longtable}{|p{3.3cm}|p{10.5cm}|}
\hline
\textbf{Model} & \textbf{Grid Searched (identical for No-PCA and With-PCA)} \\ \hline
\endhead
SVM & kernel $\in$ \{linear\} $\times$ C $\in$ \{0.1,1,10\}; kernel $\in$ \{rbf\} $\times$ C $\in$ \{0.1,1,10\} $\times$ $\gamma \in$ \{scale,auto\} \\ \hline
Naive Bayes & var\_smoothing $\in \{10^{-9}, 10^{-8}, 10^{-7}, 10^{-6}\}$ \\ \hline
KNN & n\_neighbors $\in \{3,5,7,9,11\}$, weights $\in$ \{uniform, distance\}, metric $\in$ \{euclidean, manhattan\} \\ \hline
Logistic Regression & C $\in \{0.01,0.1,1,10,100\}$, penalty=L2, solver=liblinear \\ \hline
Decision Tree & max\_depth $\in \{3,5,7,None\}$, min\_samples\_leaf $\in \{1,2,4\}$ \\ \hline
Random Forest & n\_estimators $\in \{50,100,200\}$, max\_depth $\in \{5,10,None\}$ \\ \hline
AdaBoost & n\_estimators $\in \{50,100,200\}$, learning\_rate $\in \{0.1,0.5,1.0\}$ \\ \hline
Gradient Boosting & n\_estimators $\in \{100,200\}$, learning\_rate $\in \{0.05,0.1\}$, max\_depth $\in \{2,3\}$ \\ \hline
XGBoost (HistGB) & max\_iter $\in \{100,200\}$, learning\_rate $\in \{0.05,0.1\}$, max\_depth $\in \{3,5\}$ \\ \hline
Stacking & Base models fixed (SVM, Gaussian NB, Decision Tree) $\rightarrow$ Logistic Regression meta-learner; ``hyperparameter'' here is the base-model choice itself, per the assignment's own framing \\ \hline
\end{longtable}
All grids were searched via \texttt{GridSearchCV} with 5-fold \texttt{StratifiedKFold}
(\texttt{shuffle=True, random\_state=42}), optimizing accuracy, run identically on both the
standardized original 30-feature space (No-PCA) and the 10-component PCA-reduced space
(With-PCA), with a fresh unfitted estimator instance per setting.

\section{5-Fold Cross-Validation Results (No-PCA vs. With-PCA)}

\begin{longtable}{|p{3.3cm}|p{2.3cm}|p{2.3cm}|p{2.1cm}|p{2.1cm}|}
\hline
\textbf{Model} & \textbf{Best Params (No-PCA)} & \textbf{Best Params (With-PCA)} & \textbf{Avg CV Acc (No-PCA)} & \textbf{Avg CV Acc (With-PCA)} \\ \hline
\endhead
SVM & linear, C=0.1 & linear, C=1 & 0.9758 & \textbf{0.9780} \\ \hline
Naive Bayes & smoothing=1e-9 & smoothing=1e-9 & \textbf{0.9341} & 0.9165 \\ \hline
KNN & manhattan, k=3, uniform & euclidean, k=5, uniform & \textbf{0.9714} & 0.9648 \\ \hline
Logistic Regression & C=1, L2 & C=1, L2 & 0.9802 & \textbf{0.9824} \\ \hline
Decision Tree & depth=5, leaf=4 & depth=5, leaf=1 & \textbf{0.9385} & 0.9231 \\ \hline
Random Forest & depth=10, n=100 & depth=None, n=200 & \textbf{0.9626} & 0.9604 \\ \hline
AdaBoost & lr=1.0, n=100 & lr=0.5, n=100 & \textbf{0.9802} & 0.9648 \\ \hline
Gradient Boosting & lr=0.1, depth=2, n=200 & lr=0.1, depth=3, n=200 & \textbf{0.9670} & 0.9626 \\ \hline
XGBoost (HistGB) & lr=0.05, depth=3, iter=200 & lr=0.1, depth=3, iter=100 & \textbf{0.9714} & 0.9692 \\ \hline
Stacking & SVM+NB+DT $\to$ LogReg & SVM+NB+DT $\to$ LogReg & 0.9648 & \textbf{0.9670} \\ \hline
\end{longtable}

\begin{figure}[H]
\centering
\includegraphics[width=0.95\linewidth]{no_pca_vs_pca_comparison.png}
\caption{5-fold CV accuracy comparison: No-PCA vs. With-PCA, all 10 models}
\end{figure}
\textbf{Inference:} PCA \emph{improved} average CV accuracy for exactly three models -- SVM
(+0.22pp), Logistic Regression (+0.22pp), and Stacking (+0.22pp) -- and \emph{reduced} it for
the other seven, most sharply for Naive Bayes ($-$1.76pp) and Decision Tree ($-$1.54pp). The
pattern is not random: PCA helped or was neutral for every linear/margin-based model and the
meta-learned Stacking ensemble, while it hurt every purely tree-based model (Decision Tree,
Random Forest, AdaBoost, Gradient Boosting, XGBoost) to some degree.

\begin{figure}[H]
\centering
\includegraphics[width=0.95\linewidth]{stability_comparison.png}
\caption{5-fold CV standard deviation (stability) comparison: No-PCA vs. With-PCA}
\end{figure}
\textbf{Inference:} Stability results are mixed rather than uniformly favouring either
setting. Logistic Regression, Naive Bayes, and Stacking became \emph{more} stable
(lower fold-to-fold std) with PCA, while Decision Tree, Gradient Boosting, and XGBoost
became \emph{less} stable -- for the tree-based models, this makes sense: PCA components are
linear combinations of all 30 original features, which removes the natural
axis-aligned split boundaries trees rely on, making individual tree structure (and hence
fold-to-fold variance) more sensitive to exactly which samples land in each fold.

\section{Test Set Performance (Selected Models)}
\begin{longtable}{|p{3.3cm}|p{1.6cm}|p{1.6cm}|p{1.7cm}|p{1.7cm}|}
\hline
\textbf{Model} & \textbf{Acc (No-PCA)} & \textbf{Acc (PCA)} & \textbf{F1 (No-PCA)} & \textbf{F1 (PCA)} \\ \hline
\endhead
SVM & \textbf{0.9825} & 0.9649 & \textbf{0.9861} & 0.9718 \\ \hline
Logistic Regression & \textbf{0.9825} & 0.9649 & \textbf{0.9861} & 0.9718 \\ \hline
Random Forest & \textbf{0.9561} & 0.9386 & \textbf{0.9655} & 0.9510 \\ \hline
Stacking & 0.9649 & 0.9649 & 0.9726 & 0.9722 \\ \hline
\end{longtable}
\textbf{Note:} On this single 114-row held-out test split, No-PCA edges out or ties With-PCA
for every selected model -- worth noting as a contrast with the CV table above, where SVM and
Logistic Regression's \emph{cross-validated} average was actually slightly higher with PCA.
This is the same single-split-vs-CV-average discrepancy observed in earlier experiments on
this dataset, and is why Section 6's 5-fold CV table, not this single split, is the primary
basis for the conclusions in Section 10.

\section{Result Visualization}

\begin{figure}[H]
\centering
\includegraphics[width=0.95\linewidth]{confusion_matrices_selected.png}
\caption{Confusion matrices for selected models -- top row: No-PCA, bottom row: With-PCA}
\end{figure}
\textbf{Inference:} SVM and Logistic Regression show slightly more errors under PCA than
without it on this test split (consistent with Section 7's test-accuracy dip), while
Stacking's error pattern is nearly identical in both settings -- the meta-learner appears to
compensate for whatever information the compression removes.

\begin{figure}[H]
\centering
\includegraphics[width=0.48\linewidth]{roc_curves_no_pca.png}
\includegraphics[width=0.48\linewidth]{roc_curves_with_pca.png}
\caption{ROC curves for selected models -- No-PCA (left) vs. With-PCA (right)}
\end{figure}
\textbf{Inference:} AUC values remain very high ($>$0.97) for every selected model in both
settings, and the relative ranking of models is largely preserved -- PCA does not
qualitatively change which model ranks predictions best, even where it shifts the exact
accuracy/F1 numbers at the default 0.5 threshold.

\section{Observation Questions}
\begin{enumerate}
\item \textbf{Which models improved most with PCA? Which did not? Why?} SVM, Logistic
Regression, and Stacking improved slightly (all +0.22pp CV accuracy). Naive Bayes ($-$1.76pp)
and Decision Tree ($-$1.54pp) degraded most. Linear/margin-based models improved because PCA
components are orthogonal linear combinations that can slightly reduce multicollinearity-related
noise in the decision boundary; Naive Bayes degraded because it assumes (conditional)
feature independence, and while PCA components are uncorrelated, GaussianNB's per-feature
variance estimates on rotated, mixed components no longer align with the natural boundaries of
the original measurements; Decision Tree degraded because PCA components lose the
directly-interpretable, axis-aligned thresholds (e.g. ``worst radius $\le$ 16.8'') that trees
split on naturally.
\item \textbf{Did PCA reduce variance across folds (more stable results)?} Only for some
models -- Logistic Regression, Naive Bayes, and Stacking became more stable, but Decision
Tree, Gradient Boosting, and XGBoost became \emph{less} stable (higher fold std), and SVM,
KNN, Random Forest, and AdaBoost were roughly unchanged. PCA does not uniformly stabilize
cross-validation; its effect on stability depends on how well the model family's inductive
bias matches rotated, mixed features versus original, individually meaningful ones.
\item \textbf{For high-dimensional data, was PCA beneficial in reducing overfitting?} This
30-feature dataset is only mildly high-dimensional relative to its 455 training rows, and the
results suggest PCA's overfitting-reduction benefit here was modest at best: tree-based
ensembles (which are naturally somewhat overfitting-resistant already, per Experiment 5/7's
findings) did not benefit and mostly got worse, while only the already low-variance linear
models saw a (small) improvement. PCA would likely show a clearer overfitting-reduction
benefit on datasets with a much higher feature-to-sample ratio than this one.
\item \textbf{How did linear models behave compared to ensemble models with PCA?} Linear
models (Logistic Regression, SVM) were the most PCA-friendly family in this experiment --
both improved slightly under PCA. Ensemble models (Random Forest, AdaBoost, Gradient
Boosting, XGBoost) were uniformly PCA-averse, each losing between 0.2 and 1.5 percentage
points of CV accuracy. This is a clean, consistent split along model family lines, not
scattered across both groups.
\item \textbf{Did stacking show robustness to dimensionality reduction compared to single
models?} Yes -- Stacking was one of only three models that improved under PCA (+0.22pp), and
its test-set F1 (0.9726 vs. 0.9722) was nearly unchanged between settings, the smallest
No-PCA/PCA gap of any model in Section 8's table. Combining a linear-friendly base learner
(SVM), a rotation-sensitive one (Naive Bayes), and a rotation-averse one (Decision Tree)
appears to let the meta-learner average out each base learner's individual sensitivity to the
feature transformation, producing the most dimensionality-reduction-robust result of any
model tested.
\end{enumerate}

\section{Discussion}
This experiment shows that PCA's effect is not a uniform "helps" or "hurts" verdict but
depends strongly on model family. The clean split -- linear/margin-based models and Stacking
benefiting, every tree-based model degrading -- is explainable directly from each family's
inductive bias: trees rely on axis-aligned, individually-meaningful feature thresholds that
PCA's rotation destroys, while linear models only need a good separating hyperplane, which
orthogonal, decorrelated components can provide just as well (or slightly better) than the
original correlated features. The magnitude of every effect observed here is small (all
changes under 2 percentage points), consistent with this dataset's relatively modest
dimensionality (30 features, 455 training rows) -- PCA's practical value would likely be much
larger on genuinely high-dimensional data (hundreds or thousands of features) where variance
reduction from compression matters more, or where training time reduction (fewer input
dimensions) is itself a practical goal, which was not directly measured in this experiment's
tuning-time comparisons but would be a natural extension.

\section{Conclusion}
For the Wisconsin Diagnostic Breast Cancer dataset, PCA at a 95\% variance target (compressing
30 features to 10 components) produced small, model-family-dependent effects rather than a
uniform improvement or degradation. \textbf{PCA is beneficial} for linear/margin-based models
(SVM, Logistic Regression) and for the heterogeneous Stacked Ensemble, and \textbf{PCA is not
beneficial} for tree-based models (Decision Tree, Random Forest, AdaBoost, Gradient Boosting,
XGBoost), each of which relies on original, axis-aligned feature thresholds that PCA's
rotation removes. The overall best-performing model in this experiment, SVM (linear kernel),
achieved its best cross-validated accuracy \emph{with} PCA (0.9780) even though its single
test-split accuracy was marginally higher without it -- reinforcing that hyperparameter and
preprocessing decisions should be validated by cross-validation, not a single train-test
split. The practical recommendation from this experiment is to apply PCA selectively based on
the downstream model family, rather than as a blanket preprocessing step.

\section{References}
\begin{enumerate}[label={[\arabic*]}]
\item Scikit-learn Documentation, ``PCA,'' \url{https://scikit-learn.org/stable/modules/generated/sklearn.decomposition.PCA.html}
\item Scikit-learn Documentation, ``Ensemble Methods,'' \url{https://scikit-learn.org/stable/modules/ensemble.html}
\item Scikit-learn Documentation, ``HistGradientBoostingClassifier,'' \url{https://scikit-learn.org/stable/modules/generated/sklearn.ensemble.HistGradientBoostingClassifier.html}
\item UCI Machine Learning Repository, ``Breast Cancer Wisconsin (Diagnostic) Data Set,'' \url{https://archive.ics.uci.edu}
\end{enumerate}

\end{document}
